% =============================================================================
%  UNED · Análisis Matemático II — superficie cilíndrica del restaurante
%  Ejercicio resuelto portado al sistema visual claudeeditorial.
%  Compilacion: xelatex aphorisms/uned-superficie-cilindrica.tex
% =============================================================================
\documentclass{claudeeditorial}

\renewcommand{\DocKicker}{UNED · Análisis Matemático II}
\renewcommand{\DocTitle}{Superficie cilíndrica del restaurante de la ladera}
\renewcommand{\DocSubtitle}{Cálculo del área de una pared vertical de cristal definida por $x^2+(y-R)^2=R^2$ con $0\le z\le 4R^2-2Ry$}
\renewcommand{\DocAuthor}{Patricio Hernández Ballester}
\renewcommand{\DocRole}{Profesor de apoyo · Preparador de Oposiciones}
\renewcommand{\DocDate}{24 de enero de 2025}
\renewcommand{\DocRef}{UNED-001}
\renewcommand{\DocFooter}{claudeeditorial · ejercicio resuelto}

\renewcommand{\ProjectName}{Acceso $\geq 25$ · UNED}
\renewcommand{\ProjectPhase}{Curso 2024/2025}
\renewcommand{\ProjectStatus}{Resuelto}

\hypersetup{
  pdftitle={UNED - Superficie cilindrica de la ladera},
  pdfauthor={Patricio Hernandez Ballester},
  pdfkeywords={matemáticas, superficies, cálculo de áreas, integrales, UNED}
}

\ClaudeRunningHeader

\begin{document}
\BodyCopy
\ClaudeCover

\ClaudeChapterPage{01}{Enunciado}{Construcción de un restaurante sobre la ladera de una montaña: superficie cilíndrica de cristal con altura variable según el plano del arquitecto.}

\section{Problema}

Se proyecta la construcción de un restaurante sobre la ladera de una montaña. Parte de la pared vertical está formada por la superficie cilíndrica
\[
  x^2 + (y - R)^2 = R^2,
\]
y la altura (en la variable $z$) se extiende, según el plano del arquitecto, desde $z=0$ hasta
\[
  z = 4R^2 - 2\,R\,y.
\]

Se pide determinar el \Highlight{área de la pared vertical} que se construirá con cristal, correspondiente precisamente a esa superficie (cilindro) entre $z=0$ y $z = 4R^2 - 2Ry$.

\begin{ClaudeInfo}[Objetivo]
Calcular la superficie lateral
\[
  S=\bigl\{(x,y,z)\in\R^3 : x^2+(y-R)^2=R^2,\; 0\le z\le 4R^2 -2Ry\bigr\}.
\]
\end{ClaudeInfo}

\section{Parametrización de la superficie}

Para describir el cilindro $x^2 + (y-R)^2 = R^2$ usamos coordenadas cilíndricas desplazadas en $y$:
\[
  \begin{cases}
    x(\theta,z) = R\cos\theta,\\
    y(\theta,z) = R + R\sin\theta,\\
    z(\theta,z) = z,
  \end{cases}
\]
con $\theta\in[0,2\pi]$. Como $y = R\,(1 + \sin\theta)$, el límite superior para $z$ es
\[
  z_{\max} = 4R^2 - 2R\bigl(R + R\sin\theta\bigr)
          = 2R^2\bigl(1 - \sin\theta\bigr).
\]
Por tanto, $z \in [\,0,\,2R^2(1 - \sin\theta)\,]$.

\section{Cálculo del elemento de área}

El vector de parametrización es
\[
  \vec{r}(\theta,z) = \bigl( R\cos\theta,\; R + R\sin\theta,\; z \bigr).
\]
Sus tangentes:
\[
  \frac{\partial \vec{r}}{\partial \theta}
  = (-R\sin\theta,\; R\cos\theta,\; 0),
  \qquad
  \frac{\partial \vec{r}}{\partial z}
  = (0,\;0,\;1).
\]
El producto vectorial:
\[
  \frac{\partial \vec{r}}{\partial \theta} \times \frac{\partial \vec{r}}{\partial z}
  = \bigl(R\cos\theta,\; R\sin\theta,\; 0\bigr),
\]
con norma
\[
  \norm*{\frac{\partial \vec{r}}{\partial \theta} \times \frac{\partial \vec{r}}{\partial z}}
  = \sqrt{R^2\cos^2\theta + R^2\sin^2\theta} = R.
\]

\begin{ClaudeTip}[Elemento de área]
\[
  \mathrm{d}S = R\,\mathrm{d}\theta\,\mathrm{d}z.
\]
\end{ClaudeTip}

\section{Integral de la superficie}

\[
  \iint_S 1\,\mathrm{d}S
  = \int_0^{2\pi}\!\!\int_0^{2R^2(1-\sin\theta)} R\,\mathrm{d}z\,\mathrm{d}\theta.
\]

\subsection{Evaluación}

\begin{align*}
  \int_0^{2\pi} R\,\bigl(2R^2(1-\sin\theta)\bigr)\,\mathrm{d}\theta
  &= 2R^3 \int_0^{2\pi}(1-\sin\theta)\,\mathrm{d}\theta \\
  &= 2R^3 \Bigl[\theta\bigr|_0^{2\pi} - (-\cos\theta)\bigr|_0^{2\pi}\Bigr] \\
  &= 2R^3 \bigl[2\pi - 0\bigr] = 4\pi R^3.
\end{align*}

(Se ha usado $\cos(2\pi)=\cos(0)=1$, de donde la diferencia se anula.)

\section{Conclusión}

\begin{ClaudeCheck}[Resultado]
El área de la superficie cilíndrica (la pared vertical de cristal) es
\[
  \boxed{\,4\pi R^3\,}.
\]
\end{ClaudeCheck}

Por tanto, la pared curvada vertical del restaurante tiene un área de $4\pi R^3$, expresable directamente en función del radio $R$ del cilindro generatriz.

\ClaudeSignature{\DocAuthor}{\DocRole}

\end{document}
